\documentclass{article}
\usepackage[T1]{fontenc}
\usepackage{lmodern}
\usepackage{graphicx}
\usepackage{amsmath,amssymb}
\usepackage[a4paper,margin=2cm]{geometry}
\usepackage[absolute,overlay]{textpos}
\setlength{\TPHorizModule}{1mm}
\setlength{\TPVertModule}{1mm}
\usepackage[indonesian]{babel}
\usepackage{hyperref}
\setlength{\emergencystretch}{2em}
% unit: Halaman 36

\begin{document}

% === Halaman 36 (llm) ===

\section*{Aplikasi Cipher Alir}

\begin{itemize}
    \item \textit{Cipher} alir cocok untuk mengenkripsi aliran data yang kontinu melalui saluran komunikasi, misalnya:
    \begin{enumerate}
        \item Mengenkripsi data pada saluran yang menghubungkan antara dua buah komputer.
        \item Mengenkripsi suara pada jaringan telepon \textit{mobile} GSM.
    \end{enumerate}
    \item Alasan: jika bit cipherteks yang diterima mengandung kesalahan, maka hal ini hanya menghasilkan satu bit kesalahan pada waktu dekripsi, karena tiap bit plainteks ditentukan hanya oleh satu bit cipherteks.
    \item Selain itu, alasannya karena \textit{cipher} alir itu komputasinya cepat (dibutuhkan untuk \textit{real-time processing})
    \end{itemize}

\end{document}
